\section{Collapsing the cheap parameters}
\label{c:cheap-section}

This section has two outputs.  First we move the cheap probe values to
zero while retaining a lower bound on every angular section and a bounded
form discrepancy.  We then prove a two-sided bound for all angular
values at valuations Gaussian on the cheap probes, vanishing there,
and satisfying the scaled lower bound and the stated discrepancy bound.
The bridge is the cost gap from
Proposition~\ref{b:bounded-axes-estimates}: it forces certain traced
outputs of bounded degree and cost to have exactly zero value.  The
two-sided estimate will supply the lower dimension bound in the next
section.

We continue with the actual field $K$, the fractional form $\theta$,
and the nested multiplicative angular systems $V_m$.
Degrees are nonnegative multiples of $n$, and the cost of a nonzero
$F\in V_m$, with this assigned degree, is
$b(F)=m+v(F)$.  In particular, changing the assigned degree changes the
cost.  We retain the uniform comparison $m\le Cb(F)$ for positive
costs.  Constants denoted by $C$ may increase, but are uniform along
the chosen sequence.  Let $y=(y_1,\ldots,y_r)$ be actual angular lifts
of the cheap probes in \eqref{eq:C8} and \eqref{eq:D3}.  Their degrees
and costs are uniformly bounded; their initials at $v$ are algebraically
independent.  Thus $v$ is Gaussian on $y$, and the vector $v(y)$ ranges
in a fixed bounded set.  As in the preceding section, the notation
$P_0=\mathfrak p^{\ell_0}$ refers to a Frobenius power in a reduction
of characteristic $\mathfrak p$; it is unrelated to the auxiliary
cone dimension $p$.

\par
\begingroup
  \dimen0=\pagegoal
  \advance\dimen0 by -\pagetotal
  \ifdim\dimen0<7\baselineskip
    \newpage
  \fi
\endgroup
\subsection{A translated-box count}

A section-dimension bound will control the number of distinct values of
powers of one trace output multiplied by cheap probe monomials.  The
following count turns that control into a rational linear expression in
the probe values, with bounded denominator and coefficients.  These
coefficient bounds will control the slope of the optimization objective.

\begin{lemma}\label{c:cheap-box-count}
Fix $r\ge1$ and $A>0$.  There are constants $N_0,D_0,C_0$, depending
only on $r,A$, with the following property.  Let $T\ge1$, let
$t_1,\ldots,t_r$ be linearly independent over $\Q$, and let $s\in\R$.
For an integer $N\ge N_0$, set $L=\lceil NT\rceil$.  Suppose that
\[
 \#\{as+\alpha\cdot t:0\le a<N,\quad
                    \alpha\in\{0,\ldots,L-1\}^r\}
       \le A(NT)^r.
\]
Then $s=\gamma\cdot t$ for a unique $\gamma\in\Q^r$.  The order of
$\gamma$ in $\Q^r/\Z^r$ is at most $D_0$, and
$|\gamma|\le C_0T$.
\end{lemma}

\begin{proof}
If $s$ is outside the rational span of the $t_j$, all $NL^r$ displayed
values are distinct.  This contradicts the bound once $N>A$.
Otherwise write $s=\gamma\cdot t$, uniquely by independence of $t$.
The distinct values are precisely the distinct vectors
$a\gamma+\alpha$.  If $D$ is the order of $\gamma$ modulo $\Z^r$,
these vectors occupy $\min(N,D)$ disjoint lattice cosets, each
containing a box of $L^r$ points.  Consequently
$\min(N,D)L^r\le A(NT)^r\le AL^r$; taking $N>A$ bounds $D$ by a
constant $D_0$.

Use only $a=0,D,2D,\ldots,qD$, where
$q=\lfloor(N-1)/D\rfloor$.  For a coordinate $j$, the projections of
the corresponding boxes are integer intervals of length $L$, whose
successive displacement is $D\gamma_j\in\Z$.  Their union has
\[
                 L+q\min(L,D|\gamma_j|)
\]
points.  Over each such point, one of the boxes supplies $L^{r-1}$
points in the remaining coordinates.  Hence
$q\min(L,D|\gamma_j|)\le(A-1)L$.  Choose $N_0$ so large, in terms
of $A,D_0$, that $q>A+1$ and $qD\ge N/2$.  The minimum cannot equal
$L$, and therefore
\[
             |\gamma_j|\le\frac{(A-1)L}{qD}\le C_0T.
\]
Increasing $C_0$ gives the asserted norm bound.  If $A<1$, the
hypothesis is empty, so this case causes no exception.
\end{proof}

\subsection{Moving the cheap weights to zero}

\begin{proposition}\label{c:cheap-collapse}
There are fixed sufficiently large positive constants $B_0$ and
$\kappa$ such that every sufficiently late member of the sequence
admits a monomial valuation $w_c$, Gaussian on $y$, satisfying
\[
 w_c(y)=0,\qquad A_\theta(w_c)\le\kappa,\qquad
                 w_c(F)+B_0m\ge\tfrac12 b(F)
       \quad(0\ne F\in V_m).
\]
The choices are made in the order $B_0$, then $\kappa$, before
passing to a sufficiently late member.
\end{proposition}

\begin{proof}
If $r=0$, take $w_c=v$ and $B_0\ge1$; the assertion is immediate.
Assume henceforth that $r>0$.  Use \eqref{eq:E3} with cap $q_*=1$, barrier $q_{\rm bar}=1/8$,
and budget scalar $B_*=B_0$.  Write
$L(t)=\min\mathcal L$ under the Gaussian constraint $w(y)=t$.
At $t=v(y)$ the competitor $v$ has $A_\theta(v)=0$ and $q(v)=1$
when $B_0\ge1$.  Since $A_\theta\ge0$ and $q\le1$, this gives
$L(v(y))=-\kappa$.  We shall move $t$ along the segment from $v(y)$
to zero.  Work in a fixed bounded open parameter region containing
all these segments, and consider the part where $L(t)<-\kappa/2$.
Every optimizer there has $q>1/2$.  Thus it lies strictly above
the barrier, and the openness, continuity, and local polyhedral
properties of the constrained optimization apply.

Take a full-dimensional cell, a point $t$ with rationally independent
coordinates, and the tiebroken optimizer $w$ there.  Fix its local
affine continuation by optimizers with the same active mask.
Apply Corollary~\ref{c:strict-budget} and
Proposition~\ref{c:trace-budget} on the
geometric generic fibre over the $w$-initials $g=\bar y$,
without further Taylor
shifts.  Select a homogeneous summand with surviving full-initial-field
trace and lift its factors to the actual field in a reduction.
The continuation is fixed before choosing the tiny exponents and
rational approximations to $\kappa/b_i$.  Choose them so that, after
division by $P_0$, their total errors in degree, value, and Lipschitz
variation on that continuation are bounded by a fixed small constant.
For the resulting nonzero traced output $F'\in V_{m'}$, put
\[
 \tau=\sum_i\tau_i,\qquad
 \rho=\sum_i\tau_i\frac{m_i}{b_i}\le C.
\]
Here the $\tau_i$ are the active multiplicities in \eqref{eq:E10};
$\tau\le1$, and $\tau<1$ is possible only when the cap is active.
The degree bound and the equality case of \eqref{eq:E8} give
\begin{equation}\label{eq:F1}
 \begin{split}
 m'&\le C+\kappa\rho,\\
 w(F')&=(1-P_0^{-1})
       \bigl(\kappa\tau q-\kappa B_0\rho-A_\theta(w)\bigr)+O(1).
 \end{split}
\end{equation}
Indeed, every active input satisfies
$w(F_i)=qb_i-B_0m_i$.  The Taylor flag factors from \eqref{eq:E9}
have exponents between zero and $P_0-1$ in the expanded summands.
Their normalized degrees and values, and their slopes because
$w_{\widetilde t}(y)=\widetilde t$, are bounded.  All formulas used
here are among the finite chart and filtration data retained before
reduction.  The constants in the following counting argument are
independent of the large fixed $B_0,\kappa$.

Set $T_0=C+1+\kappa\rho$, increasing the constant in
\eqref{eq:F1} if necessary.  For a sufficiently large fixed integer
$N$, consider
\[
          (F')^a y^\alpha,\qquad 0\le a<N,\quad
                          0\le\alpha_j<NT_0,\quad\alpha_j\in\Z.
\]
These products lie in a common nested degree $O(NT_0)$ and have
values at most $C'(\kappa+T_0)N$.  The same upper bound, after a
fixed enlargement, holds near $t$ on the chosen continuation.
Their number of distinct values is at most $C''(NT_0)^r$.
To justify this last assertion uniformly, before selecting the
output transfer the filtration comparisons for every degree up to
$C'''N(1+\kappa)$.  At each such degree $m$, the original
high-cost part $b>S_0$ has
\[
                    w>S_0/2-B_0m,
\]
whereas \eqref{eq:D3} bounds the dimension of its quotient by
$C_4(1+m)^r$.  For fixed $N,B_0,\kappa$, the right side of the
value inequality exceeds all the required upper cutoffs on late
members.  Retain this implication with slack by the finite-filtration
spreading argument, including the zero vector in each high-cost
subspace.  Its strict margin persists locally on the continuation.
The quotient by values above the tested cutoff therefore has the
same dimension bound in the reduction.  Distinct values below the
cutoff give independent classes in this quotient, proving the count.

Lemma~\ref{c:cheap-box-count} now yields
\[
             w(F')=\gamma\cdot t,\qquad
 \gamma\in\Q^r,\qquad
 \operatorname{ord}_{\Q^r/\Z^r}(\gamma)\le D_0,
 \qquad |\gamma|\le C'(1+\kappa\rho).
\]
The same linear expression holds on a neighborhood of $t$ on the
affine continuation.  Indeed, after fixing $F'$, the same count
applies at nearby rationally independent parameters, with the same
bounds.  There are only finitely many rational vectors of bounded
norm and denominator.  At $t$ their evaluations are distinct.
Continuity of monomial evaluation therefore excludes all but the
central vector on a sufficiently small neighborhood; density of
rationally independent parameters extends its formula to that whole
neighborhood.

Because $t$ is bounded, this gives
$w(F')\ge-C'(1+\kappa\rho)$.  In \eqref{eq:F1} use
$A_\theta(w)\ge0$ and $q,\tau\le1$.  Since $1-P_0^{-1}\ge1/2$,
we obtain
\[
             \kappa(B_0-C')\rho\le C''(1+\kappa).
\]
Thus $\rho\le C'''/B_0$ once $B_0$ and then $\kappa$ are
sufficiently large.

We next differentiate the trace inequality, using equality at the
central point rather than assuming equality throughout the
continuation.  Denote its lifted input by $G$ and its valuations by
$w_{\widetilde t}$.  The function
\[
 D(\widetilde t)=w_{\widetilde t}(F')
       -P_0^{-1}w_{\widetilde t}(G)
       +(1-P_0^{-1})A_\theta(w_{\widetilde t})
\]
is nonnegative near $t$, by \eqref{eq:E8}, and $D(t)=0$ by the
surviving initial trace.  Active entries retain
$w_{\widetilde t}(F_i)=b_iq(w_{\widetilde t})-B_0m_i$.
Consequently the variation of their normalized contribution is
$(1-P_0^{-1})\kappa\,dq$ up to the chosen small Lipschitz error.
If $\tau<1$, the cap is active and $dq=0$, so the same statement
holds.  The flag contributions and the remaining errors have a
uniformly bounded Lipschitz constant.

The optimizer value $L$ is affine on this cell and
$w_{\widetilde t}(F')=\gamma\cdot\widetilde t$.  Subtracting
$D(t)=0$ from the inequality at $t+h$ therefore gives
\[
       \gamma\cdot h+(1-P_0^{-1})\nabla L\cdot h\ge-C|h|.
\]
Using $-h$ gives the reverse bound.  Together with the bounds on
$\gamma$ and $\rho$, these inequalities imply
\begin{equation}\label{eq:F2}
                     |\nabla L|\le C_5(1+\kappa/B_0).
\end{equation}
All degree and cutoff bounds used in this deduction are uniform.
The individual reductions, Frobenius powers, and approximations may
depend on the optimizer; no uniform choice of these is required.

Let $D_y$ bound the lengths of the segments from $v(y)$ to zero.
Choose first $B_0$ so that $D_yC_5/B_0<1/4$, and then $\kappa$
so that $D_yC_5<\kappa/4$, with strict room in both inequalities.
Integrating \eqref{eq:F2} keeps $L$ strictly below $-\kappa/2$
throughout any portion of the segment reached so far.  The bound
on generic full-dimensional cells also applies along nongeneric
segments by continuity and local polyhedrality.  At a first limiting
endpoint, compactness of the relevant sublevels supplies a feasible
optimizer with the same strict objective bound.  Openness then extends the
continuation, so the whole segment is reached.  At the endpoint $q\ge1/2$ and
$A_\theta(w_c)=L(0)+\kappa q\le\kappa$, which proves the
proposition.
\end{proof}

\subsection{Exact vanishing from a cost gap}

The following elementary limiting argument will turn a bounded
initial-value output into an output of exactly zero value at the
collapsed valuation.  The conclusion is stronger than convergence
of those values to zero.
In the two-sided estimate below, a prescribed power of the tested
function will occur in every trace input.  Once the trace output has
value zero, the lower trace inequality will bound that prescribed
factor's contribution from above.

\begin{lemma}\label{c:cheap-exact-zero}
For each sequence index $i$, let $V_{m,i}$ be nested multiplicative
$k_i$-linear subspaces of a field, containing $1\in V_{0,i}$,
and let $v_i,w_i$ be real
valuations trivial on $k_i$.  Suppose there are $r$ probes $y_{j,i}$
of uniformly bounded degree, Gaussian for both valuations, such that
$|v_i(y_{j,i})|$ is uniformly bounded and $w_i(y_{j,i})=0$.
Suppose also that there are $M_i,S_i,R_i\to\infty$ and a fixed
$C_4$ with the following properties in every tested degree
$m\le M_i$.  Costs $b_i(F)=m+v_i(F)$ are nonnegative,
\[
 \begin{gathered}
 \{b_i:C_4m<b_i\le S_i\}=\varnothing,\\
 \dim_{k_i}\bigl(V_{m,i}/\{F:b_i(F)>S_i\}\bigr)
                  \le C_4(1+m)^r,\\
               b_i(F)>S_i\ \Longrightarrow\ w_i(F)>R_i.
 \end{gathered}
\]
The high-cost subspaces include zero, whose value is $+\infty$.
Then any sequence of nonzero $F_i$ of bounded degree and bounded
$v_i(F_i)$ satisfies $w_i(F_i)=0$ eventually.
\end{lemma}

\begin{proof}
Suppose that the asserted vanishing fails.  Restrict to a subsequence
on which $w_i(F_i)\ne0$, and choose a nonprincipal ultrafilter
$\mathcal U$ on it.  Let $k_{\mathcal U}=\prod_{\mathcal U}k_i$.
The nested union
\[
             A=\bigcup_m\prod_{\mathcal U}V_{m,i}
\]
is a $k_{\mathcal U}$-algebra, because the spaces are nested and
multiply with addition of degrees.  Membership in this union means
that the degree is bounded independently of $i$; it does not permit
an unbounded sequence of degrees.  Define
\[
 I=\{a=(a_i)\in A:
                   v_i(a_i)\longrightarrow_{\mathcal U}+\infty\}.
\]
In a fixed degree $m$, nonnegative cost gives $v_i(a_i)\ge-m$.
This uniform lower bound, the triangle inequality, and additivity
on products show that $I$ is an ideal.  It is proper because it
does not contain $1$.  It is prime: if neither $a$ nor $b$ belongs
to $I$, each has a finite upper valuation bound on a
$\mathcal U$-large set; their product has such a bound on the
intersection of those sets.  Thus $A/I$ is a domain.

Every $a\in I$ also satisfies
$w_i(a_i)\to_{\mathcal U}+\infty$.  This implication uses the
gap, not just slow growth of the cutoffs.  Indeed, represent $a$ in
one fixed degree $m$.  Divergence gives $b_i(a_i)>C_4m$ on a
$\mathcal U$-large set.  The gap excludes all costs between this
bound and $S_i$, hence $b_i(a_i)>S_i$ there.  The last hypothesis
then gives $w_i(a_i)>R_i\to\infty$.

The image in $A/I$ of degree $m$ has dimension at most
$C_4(1+m)^r$ over $k_{\mathcal U}$.  To see this, take more than
this many elements and representatives in $V_{m,i}$.  The finite
quotient-dimension bound gives, at each sufficiently large $i$, a
nontrivial linear combination in the high-cost subspace.  There
are finitely many coefficients; on a $\mathcal U$-large set one
fixed coefficient can be normalized to $1$.  The resulting
coefficient sequences define constants in $k_{\mathcal U}$, and
the combination has $v_i\to_{\mathcal U}\infty$.  It therefore
gives a nontrivial dependence modulo $I$.

Let $\widehat y_j$ denote the classes of the probes in $A/I$.
They are algebraically independent over $k_{\mathcal U}$.
For a nonzero polynomial with coefficients in $k_{\mathcal U}$,
choose representatives of its finitely many nonzero coefficients.
On a $\mathcal U$-large set they remain nonzero.  Gaussianity at
$v_i$ evaluates that polynomial by the minimum of finitely many
fixed monomial weights in the $v_i(y_{j,i})$.  Those weights are
uniformly bounded, so the polynomial does not belong to $I$.

It follows that $\Frac(A/I)$ is algebraic over
$k_{\mathcal U}(\widehat y_1,\ldots,\widehat y_r)$.
Otherwise some element of $A/I$ would be transcendental over that
field.  This element and the $r$ probes lie in some fixed nested
degree.  Their monomials with every exponent at most $N$ would be
linearly independent, giving $(N+1)^{r+1}$ independent elements in
a degree $O(N)$, contrary to the established $O(N^r)$ bound.
This also proves the assertion when $r=0$.

The bounded degree and bounded $v_i$-value of $F_i$ show that its
class $\widehat F$ in $A/I$ is nonzero.  Its minimal polynomial
over the cheap rational field therefore has nonzero leading and
constant coefficients.  Clearing denominators yields
\[
       \sum_{a=0}^{d_0}h_a(\widehat y)\widehat F^{\,a}=0,
 \qquad h_a\in k_{\mathcal U}[Y_1,\ldots,Y_r],
 \qquad h_0h_{d_0}\ne0.
\]
Choose representatives of all its finitely many scalar coefficients
and lift this relation to $A$.  Its residual belongs to $I$, so its
$w_i$-value tends to $+\infty$ along $\mathcal U$.
For every nonzero $h_a$, Gaussianity on $y_i$ over the whole
constant field $k_i$, together with $w_i(y_i)=0$, gives
$w_i(h_{a,i}(y_i))=0$ on a $\mathcal U$-large set.  Intersect
these finitely many sets.  If $w_i(F_i)>0$, the constant term is
the unique least term of the lifted relation, of value zero.  If
$w_i(F_i)<0$, its highest-degree term is uniquely least, of
negative value.  In either case the residual has value at most
zero, a contradiction.  No positive lower bound on
$|w_i(F_i)|$ is used.

The relation was obtained after forming the ultraproduct, but this
does not reverse the order of the choices.  Its supports, its
coefficients after denominator clearing, and all its products have
some fixed finite nested degree $m_0$.  The hypotheses were retained
for every $m\le M_i$, with $M_i\to\infty$, so they already apply
to $m_0$ on a tail.  The proof does not require spreading this newly
chosen relation.  The contradiction on any contrary subsequence
proves ordinary eventual vanishing, as asserted.
\end{proof}

\subsection{Two-sided angular budgets}

\begin{proposition}\label{c:cheap-two-sided}
Fix the constant $B_0$ of Proposition~\ref{c:cheap-collapse} and a
positive constant $c$.  For each sufficiently small fixed
$\delta>0$, consider any actual monomial valuation $w$, Gaussian
on $y$, such that $w(y)=0$, $A_\theta(w)\le1$, and
\begin{equation}\label{eq:F3}
               w(F)+\delta B_0m\ge c\delta b(F)
               \qquad(0\ne F\in V_m).
\end{equation}
For all sufficiently late sequence members, uniformly over all
these valuations and angular functions,
\begin{equation}\label{eq:F4}
                        |w(F)|\le C_6\delta b(F).
\end{equation}
The constant $C_6$ is independent of the small fixed $\delta$;
the index after which the assertion holds may depend on $\delta$.
\end{proposition}

\begin{proof}
The lower bound follows from \eqref{eq:F3} and $m\le Cb(F)$.
Degree-zero units have zero value at $w$ by centering, so they
also satisfy the assertion.  To prove the upper bound uniformly,
consider an arbitrary sequence of tests $w,F$ with $b=b(F)>0$.
We may pass to subsequences throughout.  A contradiction for every
sequence violating one uniform upper bound gives the stated
uniform conclusion.

Return to the leading construction at $v$, and let $\bar F$ be
the $v$-initial of $F$, of charge $\beta\in M_0$.  The occurring
weight comparisons give $\|\beta\|\le Cb$, and positive costs
are bounded away from zero.  By \eqref{eq:C5}, choose an actual
angular $G$ whose initial has charge $-d'\beta$, for a positive
integer $d'$, and whose degree $m_G$ and cost are at most $Cd'b$.
Thus
\[
 u=\bar F^{\,d'}\bar G\in K_1,\qquad
 \operatorname{div}_{S_1}(u)+MH_1\ge0,\qquad
 M=d'm+m_G\le Cd'b,
\]
by purity and \eqref{eq:A9}.  Here $M>0$: if $M=0$, both degrees
vanish and the opposite proportional charges have opposite
proportional costs, one strictly positive, contradicting
nonnegative degree-zero costs.

Work on a geometric generic fibre field of $K_v$ over the cheap
initial probes $\bar y$.  For a prime-divisor valuation $Q$ there,
over the geometric generic constants, restrict first to $K_v$ and
then put $P=Q|_{K_1}$.  The algebraic-extension and relative
Abhyankar argument of \eqref{eq:E4} shows that $P$ is trivial or
prime-divisorial up to scaling.  In particular, restriction
preserves Abhyankar equality, and a nontrivial subgroup of the
discrete value group remains discrete.  Moreover, $P$ is trivial
on $K_f$: the cheap residual probes contain a transcendence basis
of $K_f$, and the remaining extension is algebraic.  We have
\begin{equation}\label{eq:F5}
 \frac{P(u)}{d'}
       \le Cb\bigl(H_1(P)+s_{\exp}A_1^+(P)\bigr),
 \qquad
             A_{\theta_{v,\mathrm{rel}}}(Q)\ge A_1^+(P).
\end{equation}
The first inequality is \eqref{eq:C4} applied to the effective
pure divisor just constructed.  For the second, first apply the
relative-form comparison in \eqref{eq:E4}, then invariantize the
restriction to $K_v=K_1({\bf z}_0)$ by \eqref{eq:E1}.  The form is
\[
       \theta_v=\pm\frac{\omega_S\wedge d\log{\bf z}_0}
                              {f_h^{1/n}},
\]
so the Gauss extension of $P$ has discrepancy exactly $A_1^+(P)$.
Equality in the second inequality of \eqref{eq:F5} therefore
requires the restriction to $K_v$ itself to be Gaussian over $P$
in these torus variables.

Fix a large $\ell$ with $1/\ell\le\delta$, before taking the
sequence limit.  We now construct a strict trace budget in which
the factors $\bar F,\bar G$ are mandatory.  Give them respective
rational exponents
\[
          u_0\in[\ell/b,2\ell/b],\qquad u_0/d'.
\]
Choose a sufficiently divisible $D\in n\Z_{>0}$ for which
$DH_1$ is free Cartier.  A basis of this complete pure angular
system has minimum order $-DH_1(P)$ at every $P$, by freeness
and \eqref{eq:A9}.  Lift it by charge-zero angular functions and
give its minimum a rational exponent in
$[K_0\ell/D,2K_0\ell/D]$, where $K_0$ is a sufficiently large
fixed constant.  The mandatory contribution equals
$u_0P(u)/d'$ and is at most
$C\ell(H_1(P)+s_{\exp}A_1^+(P))$.  The pure system subtracts
at least $K_0\ell H_1(P)$.

For $P\ne0$, choose $K_0>C$ and then take a sufficiently late
member so that $C\ell s_{\exp}<1/2$.  Since $H_1(P),A_1^+(P)$
are nonnegative and their sum is positive, these contributions
have sum strictly less than $A_1^+(P)$, and hence less than
$A_{\theta_{v,\mathrm{rel}}}(Q)$.
For $P=0$, both contributions vanish.  Failure of strictness would
force equality in \eqref{eq:F5}, so the restriction to $K_v$ would
be a Gauss valuation over the trivial valuation of $K_1$.
That restriction is nontrivial, by the algebraic-extension argument
in \eqref{eq:E4}; hence some torus basis weight is nonzero.
Add finite minimum systems consisting of $1$ and angular initials
of both signs of the basis charges, allowing multiples as supplied
by \eqref{eq:C5}.  One of their minima is negative in each such
direction.  Give each an arbitrarily tiny positive rational
exponent.  Finally add systems $(1,\bar U)$ for angular initials
generating the relative field, again with arbitrarily tiny positive
rational exponents, as in \eqref{eq:E5}.  These extra minima are
nonpositive everywhere, make the budget strict also when $P=0$,
and retain birationality through their ratios.

Clear all rational denominators.  The resulting finite system $I$
and integer $k_I$ satisfy \eqref{eq:E6}.  Every member of $I$
contains the mandatory factors, while products of the finite
minimum systems realize exactly the specified remaining minimum
orders.  Because all terms have the same mandatory factor,
retaining the field-generating ratios is unaffected by that factor.
Choose the tiny exponents to make their total normalized degree
and absolute initial-value contributions bounded by a fixed
constant.  These choices may depend on the individual test.

Before invoking trace existence, also choose the following finite
data in its reductions.  Along the sequence take slow cutoffs
\[
 M_i\longrightarrow\infty,\qquad
 M_i\ll S_i'\le S_{0,i},\qquad S_i'\longrightarrow\infty,
 \qquad R_i'\longrightarrow\infty.
\]
For every degree $m\le M_i$, retain nonnegative costs, the exclusion of all costs in
$C_4m<b\le S_i'$, and the quotient-dimension bound
$C_4(1+m)^r$ from \eqref{eq:D3}.  Retain also
\[
                 b>S_i'\quad\Longrightarrow\quad w>R_i'.
\]
This is possible by \eqref{eq:F3}: its lower bound is
$c\delta S_i'-\delta B_0M_i$ on these high parts, and we choose
$R_i'$ below this bound with slack.  Here $\delta$ is fixed and
$S_i'$ dominates $M_i$.  The finite-filtration comparisons from
the preceding section retain these statements for whole spaces
over the enlarged reduced constant field, using adapted bases and
high module tails.  There are only finitely many degrees and
cutoffs at each index.  Slowing the diagonal choices as necessary
allows all of them to be retained before choosing $P_0$ or a
traced output.

We also retain Gaussianity on $y$ at both valuations over the
entire new constant field, with bounded $v(y)$ and $w(y)=0$.
This does not follow merely from preserving finitely many polynomial
tests.  On each of the two monomial charts, write the probe initials
as rational functions of the stratum coordinates and independent
normal variables.  In characteristic zero these initials are
algebraically independent.  Retain their finite expressions and a
nonzero differential minor witnessing this independence, together
with smooth chart data and separating residue coordinates.
A sufficiently general large-characteristic reduction preserves
that minor and those separating conditions.  Over the perfect
extended constants the reduced initials are therefore still
algebraically independent.  This proves the Gaussian rule for
every polynomial over those constants, including coefficients
chosen later and varying with $i$.  The actual charts at $v,w$,
the form identities, the prescribed input functions, and their
needed values are spread simultaneously with these data.

Now apply \eqref{eq:E7} and \eqref{eq:E9} at $v$.  A surviving
full-initial-field trace, followed by homogeneous expansion and
selection of a surviving term, lifts to an actual nonzero traced
output $F'$ in the reduction.  All its input degrees and absolute
$v$-values, after division by $P_0$, are at most $C\ell$.
Indeed, $m,|v(F)|\le Cb$ and
$m_G,|v(G)|\le Cd'b$; multiplication by $u_0$ and $u_0/d'$
respectively cancels these costs.  The pure system has degree $D$
and initial value zero, while its exponent times $D$ is at most
$2K_0\ell$.  The tiny systems were chosen with bounded normalized
contributions, and the Taylor flags contribute bounded amounts.
The degree bound for actual trace and the equality case of
\eqref{eq:E8}, using $A_\theta(v)=0$, therefore give
\[
                   m'\le C\ell+n,
                  \qquad |v(F')|\le C\ell.
\]
These are bounds before knowing which surviving term is selected;
no bound on the Frobenius iterate is needed.

Apply Lemma~\ref{c:cheap-exact-zero} to these reduced systems,
with $S_i=S_i'$ and $R_i=R_i'$.  Its hypotheses were retained
before trace existence, and $\ell$ is fixed, so the output degree
and initial value are uniformly bounded.  It follows that
\[
                              w(F')=0
\]
eventually.  In particular, this is an exact statement even if a
hypothetical sequence of nonzero output values were tending to zero.
The diagonal degree tests and full Gaussianity are precisely what
allow the lemma to handle the polynomial relation chosen later.

Finally apply only the lower inequality of \eqref{eq:E8} at $w$
to the product traced to this $F'$.  The mandatory input power of
$F$ is exactly $F^{(P_0-1)u_0}$, by denominator clearing and the
chosen characteristic congruence.  Every remaining factor has
value at least $-\delta B_0$ times its degree by \eqref{eq:F3}
and the retained comparisons; constants have value zero.  Their
total degree divided by $P_0$ is at most $C\ell$.  The chart
calculation also retains the value of $F$ and
$A_\theta(w)\le1$.  Hence
\[
 0=w(F')\ge(1-P_0^{-1})\bigl(u_0w(F)-1\bigr)
                                      -C\delta B_0\ell.
\]
Since $1-P_0^{-1}\ge1/2$, $u_0\ge\ell/b$, and
$1/\ell\le\delta$, this yields
\[
            w(F)\le b\bigl(\ell^{-1}+2C\delta B_0\bigr)
                        \le C_6\delta b.
\]
The constants are independent of small fixed $\delta$ and of the
chosen sequence of tests.  This proves the upper bound and
completes \eqref{eq:F4}.
\end{proof}
